\documentclass[11pt,a4paper]{article}

\usepackage[utf8]{inputenc}
\usepackage[T1]{fontenc}
\usepackage[english]{babel}
\usepackage{amsmath,amsfonts,amssymb}
\usepackage{graphicx}
\usepackage{hyperref}
\usepackage[margin=1in]{geometry}
\usepackage{cite}

\title{\textbf{Scalable Second-Order and Hybrid Methods for Optimization with Discontinuous Regularizations}\\
\large FAPESP BEPE Research Proposal}

\author{
    \textbf{Candidate:} Gabriel Belém Barbosa \\
    \textbf{Home Advisor:} Prof. Dr. Paulo José da Silva e Silva (UNICAMP) \\
    \textbf{Host Supervisor:} Prof. Dr. Dominique Orban (Polytechnique Montréal)
}
\date{Proposed period: September 2026 -- August 2027}

\begin{document}

\maketitle

\begin{abstract}
This research project proposes a research internship abroad (BEPE) at Polytechnique Montréal under the supervision of Prof. Dominique Orban. The main objective is to develop scalable second-order and hybrid methods for large-scale optimization problems involving discontinuous regularizations, particularly cardinality constraints ($L_0$ norm) and group sparsity. The proposal bridges the highly efficient first-order heuristics (Nonmonotone Spectral Proximal Gradient - NSPG, coordinate descent, and local combinatorial algorithms) developed during the PhD in Brazil with the advanced nonsmooth optimization machinery, such as trust-region and cubic regularization methods, researched by the host group. Beyond theoretical and algorithmic development, the project envisions robust implementations within the open-source \texttt{JuliaSmoothOptimizers} (JSO) ecosystem, providing direct contributions to the optimization software community.
\end{abstract}

\tableofcontents
\newpage

\section{Introduction and Justification}

The $L_0$ regularization and its generalizations for group sparsity are fundamental tools in machine learning and signal processing, enabling variable selection and the interpretability of complex mathematical models. Due to the nonconvex and combinatorial nature of the $L_0$ penalty, efficiently solving these large-scale problems remains an open challenge in optimization. Other regularizations, such as $L_1$ (Lasso) and $L_2$ (ridge) hybrids with $L_0$, are also of recent interest for their potential to avoid overfitting while maintaining the benefits of pure $L_0$ regularization \cite{fastselect}.

During the ongoing PhD research in Brazil, the central focus has been the development of highly efficient first-order heuristics for these problem classes. Specifically, we have developed algorithms based on the Nonmonotone Spectral Proximal Gradient (NSPG) combined with Coordinate Descent (CD) and combinatorial local search methods. These nonmonotone methods act efficiently as a preprocessing step (rough support identification) for state-of-the-art algorithms (e.g., PGCCD \cite{fastselect}), overcoming dimensionality challenges in high complexity and parameter correlation environments. The implementation was developed using the Julia language, which allows high performance and flexibility in developing optimization algorithms.

Despite the high computational efficiency of first-order methods (spectral), their capacity to act as global preprocessors and identify good regions of the search space could benefit substantially from second-order approaches. Enhancing this stage with structured proximal operators allows the algorithm to get much closer to the correct manifold, ensuring that the subsequent coordinate descent and combinatorial algorithms have an easier time performing local refinement and finding overall better solutions. It is in this context that the proposed international collaboration is founded.

\subsection{Justification for the Host Institution}
The research group led by Prof. Dominique Orban at Polytechnique Montréal (GERAD) possesses internationally recognized leadership and excellence in the development of scalable nonsmooth and second-order optimization methods. Recent works from the group, such as the Levenberg-Marquardt method for nonsmooth regularized least squares \cite{aravkin2024levenberg} and scalable adaptive cubic regularization methods \cite{dussault2024scalable}, represent the state-of-the-art in the robust resolution of complex optimization problems. Additionally, recent analyses on the complexity of trust-region methods for nonsmooth optimization \cite{leconte2023complexity} offer a solid theoretical foundation for extensions into discontinuous regularizations of interest. Furthermore, the group has demonstrated a strong interest in the study of nonmonotone algorithmic behavior and the development of alternative proximal operators \cite{Leconte_2024}, which closely aligns with the applicant's current research trajectory.

Beyond theoretical alignment, the host group is the main developer of \texttt{JuliaSmoothOptimizers} (JSO), a widely used ecosystem for optimization in Julia. Immersion in this development environment will foster a unique synergy, allowing not only the theoretical advancement of the research initiated in Brazil but also its conversion into impactful, globally used software. Thus, the research internship abroad will bring a substantial contribution that could not be fully obtained solely at the home institution.

\section{Objectives}

The main objective of this research internship abroad is to extend and accelerate the resolution of optimization problems with discontinuous regularizations (such as the $L_0$ penalty and group cardinality constraints) by exploring new second-order alternatives and structured proximal operators. The goal is to utilize these methods as powerful global preprocessors that can effectively navigate complex landscapes and approach the correct active manifold. This enhanced identification of good regions will subsequently allow local refiners---such as coordinate descent and combinatorial searches---to find higher-quality solutions with greater ease.

The specific objectives include:
\begin{itemize}
    \item \textbf{Hybrid Algorithmic Development:} Design a hybrid algorithmic framework where novel structured second-order alternatives and proximal operators (e.g., block-tridiagonal variable metrics) act as global identifiers of good regions. This will be integrated with the coordinate descent and combinatorial local searches studied during the PhD, which will act as the local refiners to isolate the final optimal support;
    \item \textbf{Theoretical Convergence Analysis:} Analyze the local and global convergence of the resulting hybrid method, exploring recent complexity analyses for nonsmooth optimization \cite{leconte2023complexity};
    \item \textbf{Practical Case Studies:} Validate the proposed algorithms on real-world applications of sparse machine learning and large-scale signal processing;
    \item \textbf{Software Integration:} Implement resulting algorithms in Julia, integrating the methodologies into the \texttt{JuliaSmoothOptimizers} (JSO) ecosystem if appropriate, ensuring high scalability and making the tools available to the scientific community.
\end{itemize}

\section{Methodology}

The project will be developed in incremental and iterative stages, ensuring the methodological rigor characteristic of continuous optimization.

\subsection{Directed Study and Literature Review}
Initially, an in-depth study of the trust-region methods for nonsmooth optimization \cite{leconte2023complexity}, nonmonotone algorithmic behavior, and indefinite proximal gradient methods \cite{Leconte_2024} developed by the group at Polytechnique Montréal will be conducted. In parallel, there will be focused training on the architecture of \texttt{JuliaSmoothOptimizers} to ensure that subsequent developments are aligned with the software engineering best practices of the community.

\subsection{Development of Active Hybrid Algorithms}
The research methodology will focus on the design of a hybrid algorithmic model. The component developed in Brazil (NSPG + combinatorial searches) will act in the initial iterations of the optimization process, promoting sparsity and rapidly identifying an active set (the non-zero variables).

Additionally, we will investigate the use of structured second-order matrices in the proximal operator subproblem (Variable Metric Proximal operators). While recent approaches utilizing simple indefinite diagonal metrics have been shown to underperform compared to standard spectral methods \cite{Leconte_2024}, leveraging slightly denser structures, such as tridiagonal matrices, could provide a significant convergence boost over the spectral (identity multiple) model. In fact, our preliminary experiments on L0-regularized least squares with tridiagonal structures have already shown great promise, demonstrating that these subproblems remain computationally fast while offering high precision in support recovery.

Once the active set is identified by the global preprocessors, the algorithm will leverage the coordinate descent and combinatorial searches to strictly refine the solution over the non-zero variables. This \textit{active-set} strategy drastically reduces the computational complexity, combining the rapid landscape navigation of structured proximal operators with the precise local convergence and support isolation of the CD refiners.

\subsection{Numerical Evaluation and Benchmarking}
The implementation will be done in the Julia language, due to its high performance, which aligns with the scope of the project in Brazil and the JSO infrastructure in Canada. Extensive experiments will be conducted comparing the proposed hybrid method against state-of-the-art algorithms (such as PGCCD and the L0Learn package) on high-dimensional synthetic instances and real statistical learning datasets. The evaluation metrics will include CPU time, support precision, scalability, and the final solution quality in terms of the restricted objective function.

\section{Work Plan and Schedule}

The research internship will last 12 months (September 2026 to August 2027) and will be divided into the following main phases:

\begin{enumerate}
    \item \textbf{Month 1-2: Integration and Study (Phase 1)} \\
    Settling at the host institution; familiarization with the \texttt{JuliaSmoothOptimizers} (JSO) infrastructure; study of trust-region methods and structured proximal operators.
    \item \textbf{Month 3-5: Hybrid Algorithmic Design (Phase 2)} \\
    Theoretical development of structured proximal operators as global preprocessors and their integration with the local coordinate descent and combinatorial refiners studied in Brazil. Preliminary feasibility analysis.
    \item \textbf{Month 6-8: Computational Implementation and Partial Report (Phase 3)} \\
    Parallel implementation of the designed algorithms in the Julia language, integrating the models developed in Brazil into the JSO framework. Elaboration of the annual FAPESP Scientific Report (February 2027).
    \item \textbf{Month 9-10: Benchmarking and Validation (Phase 4)} \\
    Execution of scalable test batteries on linear and logistic regression problems with $L_0$ penalization and group sparsity. Comparison against established packages in the literature.
    \item \textbf{Month 11-12: Consolidation and Dissemination (Phase 5)} \\
    Consolidation of the results, integration of the developed methodologies into the PhD thesis, and elaboration of the BEPE Final Scientific Report required by FAPESP.
\end{enumerate}

The knowledge transfer flow will take place through regular online follow-up meetings with the advisor in Brazil, ensuring that all code development and theoretical advances are directly incorporated into the PhD thesis.


\bibliographystyle{plain}
\bibliography{bibliografia}

\end{document}
